\section{Algebraic Structure}
\label{sec:topology}

\S\ref{sec:foundations} defined the combinatorial objects that populate an astrolabe file: atoms, $k$-forms, and their stage decomposition. We now introduce algebraic structure. The uniform vertex pool property (Proposition~\ref{prop:uniform-pool}) ensures that within each stage layer $\Delta^{(m)} = \mathcal{S}_{m+1} \setminus \mathcal{S}_m$, every entry draws its references from the same hash pool $\mathcal{H}_m$. This makes each layer a natural home for the classical constructions of algebraic topology: we rename $k$-forms as $k$-simplices, define restriction maps and the boundary operator, and verify $\partial^2 = 0$.

%%% §3.1 — From Forms to Simplices (pure concept, no algebra)

\subsection{From Forms to Simplices}

\begin{definition}[Stage layer and its simplices]
\label{def:stage-layer}
The \emph{stage $m$ layer} is $\Delta^{(m)} = \mathcal{S}_{m+1} \setminus \mathcal{S}_m$, with vertex set $\mathcal{H}_m$. Its \emph{$k$-simplices} are
\[
\Delta^{(m)}_k = \mathbb{A}_k \cap \Delta^{(m)},
\]
the entries of degree $k$ that belong to stage $m$. The graded decomposition is $\Delta^{(m)} = \bigsqcup_{k \geq 1} \Delta^{(m)}_k$.
\end{definition}

The passage from $k$-form (\S\ref{sec:foundations}) to $k$-simplex is a change of perspective, not of data: the same entry in \texttt{astrolabe.json} is viewed as a combinatorial object in \S\ref{sec:foundations} and as an algebraic object carrying boundary relations here.

\begin{example}[Simplices within $\Delta^{(0)}$]
\label{ex:delta0}
Consider an astrolabe file with five atoms:
\begin{verbatim}
  "a": {"ref": ["a"], "record": {"name": "Nat.add_comm"}}
  "b": {"ref": ["b"], "record": {"name": "Nat.zero_add"}}
  "c": {"ref": ["c"], "record": {"name": "Nat.succ_pred"}}
  "d": {"ref": ["d"], "record": {"name": "List.length"}}
  "e": {"ref": ["e"], "record": {"name": "List.append"}}
\end{verbatim}
The vertex set is $\mathcal{H}_0 = \{\texttt{a}, \texttt{b}, \texttt{c}, \texttt{d}, \texttt{e}\}$. The stage layer $\Delta^{(0)} = \mathcal{S}_1 \setminus \mathcal{S}_0$ contains all entries whose refs are ordered tuples from $\mathcal{H}_0$:

\medskip
\noindent\textbf{$\Delta^{(0)}_1$: 1-simplices} ($\mathbb{A}_1 \cap \Delta^{(0)}$, binary relations between atoms):
\begin{verbatim}
  "e1": {"ref": ["a", "b"], "record": {"sort": "uses"}}
  "e2": {"ref": ["b", "c"], "record": {"sort": "uses"}}
  "e3": {"ref": ["a", "c"], "record": {"sort": "generalizes"}}
\end{verbatim}
These are edges in the classical sense: \texttt{e1} records that \texttt{Nat.add\_comm} uses \texttt{Nat.zero\_add}. The ordered pair $[\texttt{a}, \texttt{b}] \neq [\texttt{b}, \texttt{a}]$: the direction matters.

\medskip
\noindent\textbf{$\Delta^{(0)}_2$: 2-simplices} ($\mathbb{A}_2 \cap \Delta^{(0)}$, ternary relations among atoms):
\begin{verbatim}
  "f1": {"ref": ["a", "b", "c"], "record": {"sort": "proof_chain"}}
\end{verbatim}
This records that \texttt{a}, \texttt{b}, \texttt{c} participate in a joint proof chain.

\medskip
All refs in $\Delta^{(0)}$ draw from $\mathcal{H}_0 = \{\texttt{a}, \texttt{b}, \texttt{c}, \texttt{d}, \texttt{e}\}$. This is a classical semi-simplicial set: vertices are atoms, simplices are ordered tuples of atoms, and $\partial$ is well-defined.
\end{example}

\begin{example}[Simplices within $\Delta^{(1)}$]
\label{ex:delta1}
Now consider entries whose refs include hashes from $\mathcal{H}_1 = \mathcal{H}_0 \cup \{\texttt{e1}, \texttt{e2}, \texttt{e3}, \texttt{f1}\}$. These form $\Delta^{(1)} = \mathcal{S}_2 \setminus \mathcal{S}_1$, with vertex set $\mathcal{H}_1$.
\begin{verbatim}
  "m1": {"ref": ["e1", "e2"],       "record": {"sort": "parallel"}}
  "m2": {"ref": ["a", "e1"],        "record": {"sort": "atom_edge_link"}}
  "m3": {"ref": ["e1", "e2", "f1"], "record": {"sort": "subsumption"}}
\end{verbatim}
In $\Delta^{(1)}$, the ``vertices'' include the edges and faces from $\Delta^{(0)}$:
\begin{itemize}
\item $\texttt{m1} \in \Delta^{(1)}_1$: a 1-simplex, a relation between two edges.
\item $\texttt{m2} \in \Delta^{(1)}_1$: a 1-simplex, a relation between an atom and an edge.
\item $\texttt{m3} \in \Delta^{(1)}_2$: a 2-simplex, a ternary relation among two edges and a face.
\end{itemize}
\end{example}

\begin{definition}[Semantic simplicial complex at stage $m$]
\label{def:semantic-complex}
The \emph{semantic simplicial complex at stage $m$} is
\[
\Delta^{(m)} = \bigsqcup_{k \geq 1} \Delta^{(m)}_k, \qquad \Delta^{(m)}_k = \bigl\{\sigma \in \Delta^{(m)} \;\big|\; |\texttt{ref}(\sigma)| = k+1\bigr\},
\]
with vertex set $\mathcal{H}_m$. Atoms in $\mathcal{H}_m$ are vertices, not simplices.
\end{definition}

\begin{table}[H]
\centering
\small
\begin{tabular}{ll}
\hline
\textbf{\S\ref{sec:foundations} (combinatorics)} & \textbf{\S\ref{sec:topology} (algebra)} \\
\hline
atom & vertex of $\Delta^{(m)}$ \\
$k$-form & $k$-simplex \\
$\mathbb{A}_k$ & $\Delta^{(m)}_k$ \\
stage layer & $\Delta^{(m)}$ \\
$\texttt{ref}$ length & dimension $k$ \\
profile $\boldsymbol{\mu}$ & $\mathbb{Z}^N$-grading of multicomplex \\
restriction map $d_i$ & face map (by position) \\
& $\partial^{(j)}$ profile boundary (by vertex) \\
& $\partial = \sum_j \partial^{(j)}$ total boundary \\
\hline
\end{tabular}
\caption{Terminology correspondence. The same entry is a form in \S\ref{sec:foundations} and a simplex in \S\ref{sec:topology}; the data does not change.}
\label{tab:form-simplex}
\end{table}

%%% §3.2 — Restriction Maps and Boundary (algebraic core)

\subsection{Restriction Maps and Boundary}

With the stage layers and their simplices in place, we introduce the algebraic operations that make each layer a chain complex.

\begin{definition}[Restriction maps]
\label{def:restriction-maps}
For $\sigma \in \Delta^{(m)}$ with $\texttt{ref} = [v_0, \ldots, v_k]$, $v_i \in \mathcal{H}_m$, the \emph{$i$-th restriction map} is
\[
d_i(\sigma) = [v_0, \ldots, \hat{v}_i, \ldots, v_k], \qquad 0 \leq i \leq k.
\]
\end{definition}

\begin{example}[Restriction maps of a 2-simplex]
Recall \texttt{f1} from Example~\ref{ex:delta0} with $\texttt{ref} = [\texttt{a}, \texttt{b}, \texttt{c}]$. Its three restriction maps are:
\begin{itemize}
\item $d_0(\texttt{f1}) = [\texttt{b}, \texttt{c}]$: the relation when \texttt{a} is removed.
\item $d_1(\texttt{f1}) = [\texttt{a}, \texttt{c}]$: the relation when \texttt{b} is removed.
\item $d_2(\texttt{f1}) = [\texttt{a}, \texttt{b}]$: the relation when \texttt{c} is removed.
\end{itemize}
Note that $d_2(\texttt{f1}) = [\texttt{a}, \texttt{b}]$ matches \texttt{e1}, and $d_0(\texttt{f1}) = [\texttt{b}, \texttt{c}]$ matches \texttt{e2}. But $d_1(\texttt{f1}) = [\texttt{a}, \texttt{c}]$ matches \texttt{e3}, which happens to exist. If \texttt{e3} did not exist, the face $[\texttt{a}, \texttt{c}]$ would be a \emph{semantic void}: a relation implied by the 2-simplex but not recorded.
\end{example}

\subsubsection*{Profile boundary operators}

The restriction map $d_i$ removes the vertex at position $i$, regardless of which vertex it is. We now group these maps by \emph{which vertex} they remove, obtaining one operator per vertex.

\begin{definition}[Profile boundary]\label{def:profile-boundary}
For $j \in \{1, \ldots, N\}$ with $\mu_j > 0$, the $j$-th profile boundary operator $\partial^{(j)}: C_{\boldsymbol{\mu}} \to C_{\boldsymbol{\mu} - \mathbf{e}_j}$ is
\[
  \partial^{(j)}(\sigma) = \sum_{\substack{i = 0 \\ v_i = h_j}}^{k} (-1)^i\, d_i(\sigma).
\]
It sums over all positions occupied by vertex $h_j$, with alternating signs. When $\mu_j = 0$, we set $\partial^{(j)} = 0$.
\end{definition}

\begin{example}[Profile boundary of a 2-simplex]\label{ex:profile-boundary}
Recall \texttt{f1} from Example~\ref{ex:delta0} with $\texttt{ref} = [\texttt{a}, \texttt{b}, \texttt{c}]$, profile $\boldsymbol{\mu} = (1, 1, 1, 0, 0)$. Each vertex appears once, so each $\partial^{(j)}$ has a single term:
\begin{align*}
  \partial^{(\texttt{a})}(\texttt{f1}) &= (-1)^0\,[\texttt{b}, \texttt{c}] = [\texttt{b}, \texttt{c}], \\
  \partial^{(\texttt{b})}(\texttt{f1}) &= (-1)^1\,[\texttt{a}, \texttt{c}] = -[\texttt{a}, \texttt{c}], \\
  \partial^{(\texttt{c})}(\texttt{f1}) &= (-1)^2\,[\texttt{a}, \texttt{b}] = [\texttt{a}, \texttt{b}].
\end{align*}
Each $\partial^{(j)}$ maps $C_{(1,1,1,0,0)}$ to $C_{(1,1,1,0,0) - \mathbf{e}_j}$: three different target profiles in the profile lattice.
\end{example}

\begin{example}[Profile boundary with repeated vertex]\label{ex:profile-boundary-repeat}
Let $\sigma = [\texttt{a}, \texttt{a}, \texttt{b}]$ with profile $\boldsymbol{\mu} = (2, 1, 0, 0, 0)$.
\begin{align*}
  \partial^{(\texttt{a})}(\sigma) &= (-1)^0\,[\texttt{a}, \texttt{b}] + (-1)^1\,[\texttt{a}, \texttt{b}]
    = [\texttt{a}, \texttt{b}] - [\texttt{a}, \texttt{b}] = 0, \\
  \partial^{(\texttt{b})}(\sigma) &= (-1)^2\,[\texttt{a}, \texttt{a}] = [\texttt{a}, \texttt{a}].
\end{align*}
The two copies of \texttt{a} contribute opposite signs and cancel. This is a profile-level phenomenon with no classical analogue: the classical boundary $\partial(\sigma) = 0 + [\texttt{a}, \texttt{a}] = [\texttt{a}, \texttt{a}]$ does not reveal that the vanishing comes entirely from the \texttt{a}-direction.
\end{example}

\begin{theorem}[Multicomplex theorem]\label{thm:multicomplex}
For all $j, l \in \{1, \ldots, N\}$:
\[
  \boxed{\partial^{(j)} \circ \partial^{(l)} + \partial^{(l)} \circ \partial^{(j)} = 0.}
\]
In particular, $(\partial^{(j)})^2 = 0$ for each $j$. The profile chain groups $\{C_{\boldsymbol{\mu}}\}_{\boldsymbol{\mu} \in \mathcal{P}_m}$ with operators $\partial^{(1)}, \ldots, \partial^{(N)}$ form a $\mathbb{Z}^N$-graded $N$-fold multicomplex.
\end{theorem}

\begin{proof}
It suffices to show that every pair of positions $(i, p)$ with $i \neq p$ contributes zero to $\partial^{(j)} \circ \partial^{(l)} + \partial^{(l)} \circ \partial^{(j)}$.

Fix $\sigma = [v_0, \ldots, v_k]$ and positions $i, p$ with $v_i = h_j$, $v_p = h_l$, $i \neq p$. Both compositions yield the same $(k{-}2)$-simplex $[v_0, \ldots, \hat{v}_{\min(i,p)}, \ldots, \hat{v}_{\max(i,p)}, \ldots, v_k]$.

\emph{Case $i < p$:} $\partial^{(l)} \circ \partial^{(j)}$: remove position $i$ (sign $(-1)^i$), then position $p{-}1$ (sign $(-1)^{p-1}$). Total: $(-1)^{i+p-1}$. $\partial^{(j)} \circ \partial^{(l)}$: remove position $p$ (sign $(-1)^p$), then position $i$ (sign $(-1)^{i}$). Total: $(-1)^{i+p}$. Sum: $(-1)^{i+p-1} + (-1)^{i+p} = 0$.

\emph{Case $i > p$:} symmetric. The case $j = l$ (both positions are $h_j$) follows identically.
\end{proof}

The multicomplex theorem is the fundamental algebraic fact of the profile decomposition. All subsequent structure flows from it.

\subsubsection*{Classical operators as projections}

\begin{definition}[Total boundary operator]\label{def:total-boundary}
\label{def:boundary}
The \emph{total boundary operator} is the sum over all vertex directions:
\[
  \partial_k = \sum_{j=1}^{N} \partial^{(j)}: C_k \to C_{k-1},
  \qquad C_k = \bigoplus_{|\boldsymbol{\mu}| = k+1} C_{\boldsymbol{\mu}}.
\]
Expanding: $\partial_k(\sigma) = \sum_{j} \sum_{i:\,v_i = h_j} (-1)^i\, d_i(\sigma) = \sum_{i=0}^{k} (-1)^i\, d_i(\sigma)$, which recovers the classical formula.
\end{definition}

\begin{corollary}[$\partial^2 = 0$]\label{thm:boundary-squared}
$\partial_{k-1} \circ \partial_k = 0$ for all $k$.
\end{corollary}

\begin{proof}
$\partial^2 = \bigl(\sum_j \partial^{(j)}\bigr)^2 = \sum_j (\partial^{(j)})^2 + \sum_{j < l}(\partial^{(j)}\partial^{(l)} + \partial^{(l)}\partial^{(j)}) = 0$ by Theorem~\ref{thm:multicomplex}.
\end{proof}

\begin{corollary}[Restriction identity]\label{prop:restriction-identity}
$d_i \circ d_j = d_{j-1} \circ d_i$ for $i < j$.
\end{corollary}

\begin{proof}
This is the position-level statement underlying the proof of Theorem~\ref{thm:multicomplex}: removing $v_j$ then $v_i$ ($i < j$) is the same as removing $v_i$ then $v_{j-1}$.
\end{proof}

\begin{remark}[Logical dependence]\label{rem:logical-dependence}
In the classical presentation, one proves the restriction identity, defines $\partial = \sum (-1)^i d_i$, and deduces $\partial^2 = 0$. In the profile presentation, the logical order is reversed: one defines $\partial^{(j)}$, proves anticommutativity (Theorem~\ref{thm:multicomplex}), and obtains both $\partial^2 = 0$ and the restriction identity as corollaries. The multicomplex theorem is strictly stronger than $\partial^2 = 0$: it gives $\binom{N+1}{2}$ equations, of which $\partial^2 = 0$ is a single linear combination.
\end{remark}

\begin{remark}[Profile-level algebraic topology]\label{rem:profile-future}
The multicomplex structure of Theorem~\ref{thm:multicomplex} supports a richer algebraic topology than we develop here. Since $(\partial^{(j)})^2 = 0$ for each $j$, one can define \emph{$j$-homology} $H^{(j)}_{\boldsymbol{\mu}} = \ker \partial^{(j)} / \operatorname{im} \partial^{(j)}$ at each profile $\boldsymbol{\mu}$, yielding multigraded Betti numbers $\beta^{(j)}_{\boldsymbol{\mu}}$ that refine $\beta_k$. The anticommutativity $\partial^{(j)}\partial^{(l)} + \partial^{(l)}\partial^{(j)} = 0$ implies that $\partial^{(l)}$ induces maps on $H^{(j)}$, giving iterated homology $H^{(l,j)}$, Koszul homology $\mathcal{K}H^S$ for subsets $S \subseteq \{1, \ldots, N\}$, and a profile spectral sequence. These structures detect phenomena invisible to the total boundary $\partial$; for instance, whether two vertices have exchangeable topological roles ($H^{(l,j)} \cong H^{(j,l)}$) or are irreducibly entangled ($\mathcal{K}H^{\{j,l\}} \neq 0$). We defer the full development to a companion paper.
\end{remark}

\begin{remark}[Three-layer orientation consistency]
\label{rem:orientation}
Each 1-simplex carries three orientations that must be compatible:
\begin{enumerate}
\item \emph{Algebraic}: the ordering of $\texttt{ref} = [v_0, v_1]$ determines the sign in $\partial$.
\item \emph{Semantic}: the $\texttt{record}$ content (e.g., ``$A$ uses $B$'' vs.\ ``$B$ uses $A$'').
\item \emph{Transitive}: if $A \to B$ and $B \to C$, then the composite $A \to C$ must respect both orientations.
\end{enumerate}
Inconsistency between these layers signals a data quality issue in the knowledge base.
\end{remark}

\begin{example}[Boundary of a 1-simplex]
Let $\sigma \in \Delta^{(m)}_1$ with $\texttt{ref} = [\texttt{a1b2c3},\; \texttt{b2c3d4}]$ (an edge from \texttt{Nat.add\_comm} to \texttt{Nat.zero\_add}). Then
\[
\partial_1(\sigma) = d_0(\sigma) - d_1(\sigma) = [\texttt{b2c3d4}] - [\texttt{a1b2c3}].
\]
The boundary is ``arrive at the target, depart from the source.'' A chain of edges along a path has boundary equal to the difference of the two endpoints: all intermediate vertices cancel.
\end{example}

\begin{example}[Gene regulatory network]
In a biological knowledge base, a 1-simplex records that gene \texttt{TP53} regulates gene \texttt{MDM2}: $\texttt{ref} = [\texttt{TP53}, \texttt{MDM2}]$, with $\texttt{record} = \{(\texttt{sort}, \texttt{activates}), (\texttt{evidence}, \texttt{ChIP-seq})\}$. The boundary $\partial_1(e) = [\texttt{MDM2}] - [\texttt{TP53}]$ represents the ``flow of regulation.'' If we also have $e'$ with $\texttt{ref} = [\texttt{MDM2}, \texttt{TP53}]$ (MDM2 inhibits TP53), the chain $e + e'$ forms a cycle: $\partial_1(e + e') = 0$. This is the well-known TP53--MDM2 negative feedback loop, detected algebraically as a 1-cycle.
\end{example}

\begin{example}[Verification of $\partial^2 = 0$]
Let $\sigma \in \Delta^{(m)}_2$ with $\texttt{ref} = [\texttt{a}, \texttt{b}, \texttt{c}]$ (abbreviating hashes). Then
\[
\partial_2(\sigma) = [\texttt{b}, \texttt{c}] - [\texttt{a}, \texttt{c}] + [\texttt{a}, \texttt{b}].
\]
Applying $\partial_1$ to each term:
\[
\partial_1([\texttt{b},\texttt{c}]) = \texttt{c} - \texttt{b}, \quad
\partial_1([\texttt{a},\texttt{c}]) = \texttt{c} - \texttt{a}, \quad
\partial_1([\texttt{a},\texttt{b}]) = \texttt{b} - \texttt{a}.
\]
Therefore $\partial_1 \circ \partial_2(\sigma) = (\texttt{c} - \texttt{b}) - (\texttt{c} - \texttt{a}) + (\texttt{b} - \texttt{a}) = 0$. \checkmark

This is the total boundary; the profile decomposition (Theorem~\ref{thm:multicomplex}) gives the finer statement that each pair of vertex directions cancels independently.
\end{example}

The chain groups and boundary operators form a \emph{chain complex} under the total boundary $\partial = \sum_j \partial^{(j)}$. The profile multicomplex refines this into a $\mathbb{Z}^N$-graded diagram on the profile lattice $\mathcal{P}_m$.

\begin{figure}[ht]
\centering
\begin{tikzcd}[column sep=large]
\cdots \arrow[r, "\partial_{k+1}"] & C_k(\Delta) \arrow[r, "\partial_k"] & C_{k-1}(\Delta) \arrow[r, "\partial_{k-1}"] & \cdots \arrow[r, "\partial_2"] & C_1(\Delta) \arrow[r, "\partial_1"] & C_0(\Delta) \arrow[r] & 0
\end{tikzcd}
\caption{The chain complex of $\Delta^{(m)}$. Each arrow satisfies $\operatorname{im} \partial_{k+1} \subseteq \ker \partial_k$, and the quotient $\ker \partial_k / \operatorname{im} \partial_{k+1}$ is the $k$-th homology group.}
\label{fig:chain-complex}
\end{figure}

The condition $\operatorname{im} \partial_{k+1} \subseteq \ker \partial_k$ means: every boundary is a cycle. The quotient is the homology group.

%%% §3.3 — Homology (H_k, Betti, invariance, Euler)

\subsection{Homology}

\begin{definition}[Cycles and boundaries]
\label{def:cycles-boundaries}
The \emph{$k$-cycles} and \emph{$k$-boundaries} are the subgroups
\[
Z_k = \ker \partial_k \subseteq C_k, \qquad B_k = \operatorname{im} \partial_{k+1} \subseteq C_k.
\]
Since $\partial \circ \partial = 0$, we have $B_k \subseteq Z_k$.
\end{definition}

A \emph{cycle} is a chain of references that ``closes up'': every record entered via a reference is also exited. A \emph{boundary} is a cycle that arises as the border of a higher-dimensional structure and is therefore ``trivially closed.''

\begin{example}[Cycle vs.\ boundary in a legal code]
Three statutes $A$, $B$, $C$ reference each other cyclically via 1-simplices $e_{AB}, e_{BC}, e_{CA} \in \Delta_1$. The chain $e_{AB} + e_{BC} + e_{CA}$ is a 1-cycle ($\partial_1 = 0$). Now suppose a commentary article $f \in \Delta_2$ with $\texttt{ref} = [A, B, C]$ simultaneously references all three statutes. Then $\partial_2(f) = e_{BC} - e_{AC} + e_{AB}$ (up to sign): the cycle is a boundary. The circular cross-references are ``explained'' by the commentary, so the cycle is trivial in $H_1$. Without the commentary, the cycle represents a genuine circularity in the legal code.
\end{example}

The interesting objects are cycles that are \emph{not} boundaries: these represent genuine topological features of the knowledge base.

\begin{definition}[Homology]
\label{def:homology}
The \emph{$k$-th homology group} of $\Delta^{(m)}$ is the quotient
\[
H_k(\Delta^{(m)}) = Z_k / B_k = \ker \partial_k / \operatorname{im} \partial_{k+1}.
\]
Its rank $\beta_k = \operatorname{rank} H_k$ is the \emph{$k$-th Betti number}.
\end{definition}

The Betti numbers have direct knowledge-management interpretations. The algebraic definitions ($H_k = \ker \partial_k / \operatorname{im} \partial_{k+1}$) hold for any Astrolabe; the geometric meanings below assume the classical constraint (Definition~\ref{def:classical}):

\begin{itemize}
\item $\beta_0$: the number of \emph{connected components}: independent knowledge domains with no references between them.
\item $\beta_1$: the number of independent \emph{1-cycles}: circular chains of references that do not bound a 2-simplex. In Lean-formalized mathematics, where every dependency chain is acyclic, $\beta_1 = 0$.
\item $\beta_2$: the number of independent \emph{2-cycles}: ternary relationships that form a closed surface without being filled by a 3-simplex. These represent irreducible higher-order synergies.
\end{itemize}

\begin{example}[Circular dependency detection]
Consider three Python modules $A, B, C \in \Delta_0$ and three edges $e_{AB}, e_{BC}, e_{CA} \in \Delta_1$ with $\texttt{ref} = [A, B]$, $[B, C]$, $[C, A]$ respectively, representing \texttt{import} statements. The chain $e_{AB} + e_{BC} + e_{CA}$ is a 1-cycle ($\partial_1 = 0$: every module that is entered is also exited). If no 2-simplex fills this triangle (i.e., no record in $S$ has $\texttt{ref} = [A, B, C]$), then this cycle is not a boundary, and $\beta_1 \geq 1$. The homology class $[e_{AB} + e_{BC} + e_{CA}] \in H_1$ is a formal certificate of the circular dependency; it persists until the cycle is broken or filled.
\end{example}

\begin{example}[Knowledge island detection]
A mathematical knowledge base with $\beta_0 = 3$ has three connected components: for instance, an algebra cluster, an analysis cluster, and a topology cluster with no cross-references. Adding a single edge between any two clusters reduces $\beta_0$ by one. The Betti number $\beta_0$ quantifies the fragmentation of the knowledge base.
\end{example}

The following theorems establish that homological invariants are robust under structural operations.

\begin{definition}[Enrichment]
\label{def:enrichment}
An \emph{enrichment} is a function $E : R \to R$ that modifies only auxiliary keys: $E(r)|_{K_{\mathrm{ess}}} = r|_{K_{\mathrm{ess}}}$ and $h_{\mathrm{id}}(E(r)) = h_{\mathrm{id}}(r)$. Enrichments compose: $E_2 \circ E_1$ is again an enrichment, and enrichments writing to disjoint keys commute.
\end{definition}

\begin{theorem}[Enrichment invariance]
\label{thm:enrichment-invariance}
Let $E : S \to S$ be an enrichment: $E$ modifies only auxiliary keys and preserves $h_{\mathrm{id}}$. Then $H_k(\Delta(E(\Sigma))) \cong H_k(\Delta(\Sigma))$ for all $k$. \textup{[Lean~4]}
\end{theorem}

Adding metadata (timestamps, pagerank scores, community labels) to records does not change the homology. The topological structure depends only on the essential content and the reference structure, both of which enrichment leaves invariant.

\begin{theorem}[Faithful import preserves homology]
\label{thm:faithful-import}
Let $F : \Sigma_1 \to \Sigma_2$ be a faithful import (an injective map on simplices that preserves $\texttt{ref}$ structure). Then $F_* : H_k(\Delta(\Sigma_1)) \to H_k(\Delta(\Sigma_2))$ is injective for all $k$. \textup{[Lean~4]}
\end{theorem}

Importing a knowledge base into a larger one without modifying existing records preserves all topological features. No cycles are collapsed and no holes are filled by the import alone.

\begin{theorem}[Quotient dimension bound]
\label{thm:quotient-dimension}
Let $\Delta' = \Delta / {\sim}$ be a quotient of the Astrolabe by an equivalence relation on vertices. Then the homological dimension of $\Delta'$ does not exceed that of $\Delta$:
\[
\max\{ k \mid H_k(\Delta') \neq 0 \} \leq \max\{ k \mid H_k(\Delta) \neq 0 \}.
\]
\end{theorem}

Collapsing vertices (e.g., merging duplicate records, abstracting implementation details) can only simplify the topology, never create new higher-dimensional holes.

\begin{definition}[Euler characteristic]
\label{def:euler}
The \emph{Euler characteristic} of $\Delta^{(m)}$ is the alternating sum of Betti numbers:
\[
\chi(\Delta^{(m)}) = \sum_{k \geq 0} (-1)^k \beta_k = \sum_{k \geq 0} (-1)^k |\Delta^{(m)}_k|.
\]
\end{definition}

\begin{example}[Euler characteristic as complexity measure]
A Lean~4 library with $|\Delta_0| = 100$ declarations (vertices), $|\Delta_1| = 250$ dependencies (edges), and $|\Delta_2| = 40$ ternary interactions (faces) has $\chi = 100 - 250 + 40 = -110$. A large negative $\chi$ indicates a densely interconnected knowledge base with many cyclic structures. Refactoring that eliminates circular dependencies increases $\chi$ toward zero.
\end{example}

\begin{theorem}
$\chi(\Delta) = \chi(\Delta^{\mathrm{op}})$, where $\Delta^{\mathrm{op}}$ reverses all references (i.e., reverses the order of every $\texttt{ref}$ list).
\end{theorem}

% [MOVED TO §4] Cohomology: cochain complex, coboundary, cohomology groups,
%   H^0/H^1/H^2 interpretations, code review cost example, irreducible ternary example,
%   Kronecker pairing, cup product, cohomology ring.
%   See sections/cohomology-moved.tex

%%% §3.4 — Long Exact Sequences and Diagram Chasing

\subsection{Long Exact Sequences and Diagram Chasing}

A subcomplex $\mathcal{B} \hookrightarrow \Delta$ (a subset of records closed under taking references) induces a short exact sequence of chain complexes $0 \to C_*(\mathcal{B}) \to C_*(\Delta) \to C_*(\Delta/\mathcal{B}) \to 0$, which gives rise to the \emph{long exact sequence in homology}:

\begin{figure}[ht]
\centering
\begin{tikzcd}[column sep=small]
\cdots \arrow[r] & H_k(\mathcal{B}) \arrow[r] & H_k(\Delta) \arrow[r] & H_k(\Delta/\mathcal{B}) \arrow[r, "\delta"] & H_{k-1}(\mathcal{B}) \arrow[r] & \cdots
\end{tikzcd}
\caption{Long exact sequence. The connecting homomorphism $\delta$ detects how topological features of the quotient $\Delta/\mathcal{B}$ affect the subcomplex $\mathcal{B}$.}
\label{fig:les}
\end{figure}

\begin{example}[Microservices]
Let $\mathcal{B}$ be the subcomplex of records internal to a single service and $\Delta$ the full system. The connecting map $\delta$ identifies where circular dependencies in the full system ($H_1(\Delta/\mathcal{B}) \neq 0$) cut into the internal structure of the service ($H_0(\mathcal{B})$ splits).
\end{example}

\paragraph{Mayer--Vietoris.} If $\Delta = \mathcal{U} \cup \mathcal{V}$ is a union of two subcomplexes, the Mayer--Vietoris sequence relates their homologies:

\begin{figure}[ht]
\centering
\begin{tikzcd}[column sep=small]
\cdots \arrow[r] & H_k(\mathcal{U} \cap \mathcal{V}) \arrow[r] & H_k(\mathcal{U}) \oplus H_k(\mathcal{V}) \arrow[r] & H_k(\Delta) \arrow[r, "\delta"] & H_{k-1}(\mathcal{U} \cap \mathcal{V}) \arrow[r] & \cdots
\end{tikzcd}
\caption{Mayer--Vietoris sequence. The overlap $\mathcal{U} \cap \mathcal{V}$ controls how the homologies of the parts combine.}
\label{fig:mayer-vietoris}
\end{figure}

\begin{example}[Frontend/backend]
Let $\mathcal{U}$ = frontend records, $\mathcal{V}$ = backend records, $\mathcal{U} \cap \mathcal{V}$ = API layer. If $\delta \neq 0$, then a cycle in the full system forces the API layer to carry nontrivial topology: circular dependencies tear apart the interface.
\end{example}

\paragraph{Snake lemma.} A version upgrade $v_1 \to v_2$ induces a morphism of chain complexes, and the snake lemma produces:
\[
0 \to \ker f \to \ker g \to \ker h \xrightarrow{\delta} \operatorname{coker} f \to \operatorname{coker} g \to \operatorname{coker} h \to 0.
\]
The connecting map $\delta : \ker h \to \operatorname{coker} f$ means: deletion of inter-module relations forces creation of new intra-module structure.

\begin{example}[Python refactoring]
Deleting an import between two modules ($\ker h$) forces a new local implementation within the importing module ($\operatorname{coker} f$).
\end{example}

%%% §3.5 — Persistent Homology

\subsection{Persistent Homology}

\begin{definition}[Filtration]
\label{def:filtration}
A \emph{filtration} of $\Delta(\Sigma)$ is a nested sequence of subcomplexes
\[
\varnothing = \Delta^{(0)} \subseteq \Delta^{(1)} \subseteq \cdots \subseteq \Delta^{(n)} = \Delta(\Sigma),
\]
indexed by a parameter $\alpha_1 < \cdots < \alpha_n$. Natural filtrations include:
\begin{itemize}
\item \emph{Temporal filtration}: $\Delta^{(i)}$ contains records created before time $\alpha_i$.
\item \emph{Sort filtration}: $\Delta^{(i)}$ contains records whose sort belongs to the first $i$ levels of a type hierarchy.
\item \emph{Metric filtration}: $\Delta^{(i)}$ contains simplices with $\operatorname{vol}_k(\sigma) \geq \alpha_i$ (requires the semantic embedding of \S\ref{sec:semantics}).
\item \emph{Skeleton filtration}: $\Delta^{(i)}$ contains simplices of dimension $\leq i$.
\item \emph{Stage filtration}: $\Delta^{(i)} = \bigcup_{m \leq i} \Delta^{(m)}$, the cumulative stage decomposition.
\item \emph{Multiplicity filtration}: $\Delta^{(m)}_{\mathbf{m}}$ for increasing multiplicity bound $\mathbf{m}$ (Definition~\ref{def:max-mult}). This is a multi-parameter filtration intrinsic to the profile decomposition.
\end{itemize}
\end{definition}

\begin{example}[Temporal filtration of Mathlib]
Consider the Lean~4 mathematical library Mathlib, filtered by commit date. At $\alpha_1$ (early commits), the complex contains basic definitions and has $\beta_0 = 12$ (twelve isolated theory clusters). As lemmas connecting algebra and analysis are added ($\alpha_5$), $\beta_0$ drops to $3$. When a circular proof dependency is introduced at $\alpha_8$ and later resolved at $\alpha_{11}$, a class in $H_1$ is born at $\alpha_8$ and dies at $\alpha_{11}$, with persistence $3$. Long-lived classes indicate persistent structural features; short-lived classes indicate transient issues.
\end{example}

\begin{definition}[Persistent homology]
\label{def:persistent-homology}
The inclusion $\Delta^{(i)} \hookrightarrow \Delta^{(j)}$ for $i \leq j$ induces a map $H_k(\Delta^{(i)}) \to H_k(\Delta^{(j)})$. A homology class $\gamma \in H_k(\Delta^{(i)})$ is \emph{born} at index $i$ if $\gamma \notin \operatorname{im}(H_k(\Delta^{(i-1)}) \to H_k(\Delta^{(i)}))$, and \emph{dies} at index $j$ if its image in $H_k(\Delta^{(j)})$ becomes zero. The \emph{persistence} of $\gamma$ is $j - i$.
\end{definition}

\begin{example}[Biology: protein interaction network evolution]
A protein interaction database filtered by experimental confidence score $\alpha$ (from high-confidence to low-confidence interactions). At high $\alpha$, only well-established interactions appear, and $\beta_0$ is large (many isolated protein clusters). As $\alpha$ decreases, speculative interactions bridge clusters, reducing $\beta_0$. A cluster merger that persists across a wide range of $\alpha$ represents a robust biological pathway; one that appears only at low $\alpha$ may be an artifact.
\end{example}

\begin{remark}[Comparison with classical TDA]
Classical topological data analysis constructs simplicial complexes from geometric data (Rips complex, \v{C}ech complex) by thresholding pairwise distances. Here the complex $\Delta(\Sigma)$ is defined combinatorially by hash references, and multiple filtrations coexist on the same complex. The topology exists prior to any geometry; equipping $\Delta(\Sigma)$ with a semantic metric (\S\ref{sec:semantics}) adds a geometric filtration as one option among many.
\end{remark}
